\documentclass[10pt]{article}
\usepackage[letterpaper,margin=0.92in]{geometry}
\usepackage[T1]{fontenc}
\usepackage{lmodern,microtype,amsmath,amssymb,amsthm,mathtools}
\usepackage{graphicx,booktabs,array,natbib,enumitem,fancyhdr}
\usepackage[hidelinks]{hyperref}
\usepackage[font=small,labelfont=bf]{caption}
\setlength{\parindent}{0pt}
\setlength{\parskip}{5pt}
\setlength{\headheight}{14pt}
\pagestyle{fancy}\fancyhf{}
\fancyhead[L]{\small Finding Is Not Certifying}
\fancyhead[R]{\small Research manuscript v0.1}
\fancyfoot[C]{\thepage}
\newtheorem{theorem}{Theorem}
\newtheorem{proposition}[theorem]{Proposition}
\newtheorem{corollary}[theorem]{Corollary}
\newtheorem{lemma}[theorem]{Lemma}
\theoremstyle{definition}\newtheorem{definition}[theorem]{Definition}
\newcommand{\HH}{\mathcal H}
\newcommand{\KK}{\mathcal K}
\newcommand{\NN}{\mathcal N}
\newcommand{\UU}{\mathcal U}
\newcommand{\VV}{\mathcal V}
\newcommand{\II}{\mathcal I}
\newcommand{\FF}{\mathcal F}
\newcommand{\ind}{\mathbf 1}
\newcommand{\Cert}{\operatorname{Cert}}
\newcommand{\rank}{\operatorname{rank}}
\newcommand{\E}{\mathbb E}
\newcommand{\Prb}{\mathbb P}
\title{\vspace{-1.2em}Finding Is Not Certifying:\\Query-Limited Audits of Redundant Mechanisms}
\author{Anonymous research manuscript}
\date{15 September 2026}
\begin{document}
\maketitle
\begin{abstract}
Single-component interventions can miss joint dependencies, but finding a hidden
dependency and establishing that no others remain are different problems. We study
complete recovery of minimal harmful ablation sets for a fixed network, intervention
policy, and behavioral threshold. Under monotonicity and a declared interaction-order
bound, we characterize the smallest intervention radius permitting identification and
the exact size of a truthful identifying certificate. Building on classical independent
covers, we analyze a coverage-directed audit whose query cost is bounded by a harmonic
factor times the optimal negative certificate, plus an additive edge-extraction cost.
The argument spans all discovery phases and allows arbitrary overlap. We also show
that a network with only disjoint pair failures can require exponentially many
negative tests when higher-order alternatives must be excluded, despite admitting
logarithmic-size negative certificates under a trusted rank-two assumption.
Exhaustive finite-hypothesis checks, 2,880 synthetic audits, and audits of 20 small
trained networks validate the statements and expose their limits. Coverage-directed
queries improve the average synthetic query count over a matched random scheduler,
but do not uniformly improve neural-network auditing. Explicit cancellation and
rank-misspecification controls show why transcript consistency alone cannot justify
the assumptions behind a completeness claim.
\end{abstract}

\section{Introduction}
A model can retain its behavior after removing either of two computational pathways,
yet change substantially when both are removed. An importance score computed from
single ablations need not reveal this joint dependency. Conditional co-ablation already
uses interventions to recover compensating mechanisms in transformers
\citep{gong2026}. The outstanding question we study is narrower than whether redundancy
exists: \emph{what intervention evidence licenses a claim that the relevant joint
dependencies have all been found?}

Finding one sufficient circuit, identifying individually influential components, and
recovering a complete family of joint failure sets are different objectives. If
pathways $\{1,2\}$ and $\{3,4\}$ are individually sufficient, their four cross-pathway
pairs are minimal harmful ablations. A pruning method returning either pathway may
be correct for its own objective while not characterizing all these dependencies.
We therefore evaluate recovery of a thresholded intervention-response function,
not the size of one retained subnetwork.

The mathematical connection to hidden-hypergraph learning is established, not new.
In particular, \citet{angluin2006} explain why negative queries must form an independent
cover and how positive group tests can reveal new edges. Our contribution is an
instance-relative audit analysis that keeps three quantities separate: the intervention
ceiling, the order of alternatives excluded by the certificate, and the number of
measurements actually required. The resulting guarantees are conditional statements
about a specified intervention protocol, not identification of a unique semantic
algorithm inside a network.

We provide an exact identifiability-radius characterization and a certificate-complexity
identity. We then prove a global harmonic bound for an adaptive coverage-directed
scheduler, without multiplying the negative-certificate cost by the number of
discovered edges. A disjoint-pair construction isolates the potentially exponential
cost of excluding higher-order alternatives. Finally, our small, fully enumerable
experiments distinguish discovering all true edges from knowing that this has happened,
and test failures of the assumptions rather than silently treating them as universal.

\section{Problem and scope}
Fix a frozen model $f$, a distribution $P$, $n$ intervention units, and a replacement
policy. For $A\subseteq[n]$, let $f^{(-A)}$ denote the intervened model and define
\begin{equation}
 D(A)=\E_{x\sim P}\big[\ell(f(x),f^{(-A)}(x))\big],\qquad
 B(A)=\ind\{D(A)>\tau\},\quad \ell\in[0,1].
 \label{eq:behavior}
\end{equation}
Here ``failure'' means a specified behavioral deviation from $f$, not necessarily a
reduction in task accuracy. We require $B(\varnothing)=0$. Our initial oracle returns
exact $B(A)$; Section~\ref{sec:noise} treats sampled losses.

\textbf{Structural assumptions.} We assume monotonicity,
$A\subseteq A'\Rightarrow B(A)\le B(A')$. Let $\HH$ be the antichain of inclusion-minimal
positive masks. Then
\begin{equation}
 B_{\HH}(A)=\bigvee_{e\in\HH}\ind\{e\subseteq A\}.
 \label{eq:hypergraph}
\end{equation}
The declared hypothesis class contains \emph{all} such functions with
$\rank(\HH)=\max_{e\in\HH}|e|\le r$, with rank zero for the empty family. There is no
supplied edge-count bound: $s=|\HH|$ is a property of the target, not an input to the
audit. A known tight edge-count bound changes our lower-bound class.

Queries obey $|A|\le b$. We count distinct nonempty masks separately from the number
of input evaluations needed to label each mask. Small $b$ is a cardinality constraint;
it does not establish that interventions are on-manifold. Monotonicity and rank are
not inferred from successful reconstruction or a clean accuracy score. Lower bounds
apply to intervention-oracle evidence, not arbitrary weight inspection.

\textbf{Candidate and certificate objects.} Write
\begin{equation}
 \VV_r=\{T\subseteq[n]:1\le |T|\le r\},\qquad
 \II_r(\HH)=\{T\in\VV_r:B_{\HH}(T)=0\},\quad M=|\II_r(\HH)|.
\end{equation}
Let $\alpha(\HH)=\max\{|A|:B_{\HH}(A)=0\}$ be the largest harmless intervention size.
Define $C_{b,r}(\HH)$ as the minimum number of negative masks of size at most $b$
covering every member of $\II_r(\HH)$ by containment. Set $C=0$ if $M=0$ and $C=\infty$
if no such cover exists. This is a cardinality-constrained version of the independent
cover of \citet{angluin2006}. These quantities are used for analysis and evaluation;
the algorithm is not supplied with the true hypergraph or its optimal cover.

\section{What can the intervention domain identify?}
\begin{theorem}[Sharp identifiability radius]\label{thm:radius}
For a target $\HH$ of rank at most $r$, its full response table on masks of size at
most $b$ uniquely identifies it within the declared hypothesis class if and only if
\begin{equation}
 b\ \ge\ b_*(\HH;r):=\max\{\rank(\HH),\min(r,\alpha(\HH))\}.
 \label{eq:radius}
\end{equation}
\end{theorem}
\begin{proof}
If a minimal edge $e$ has $|e|>b$, removing $e$ from the target changes no allowed
response. If $r>b$ and $\alpha(\HH)>b$, choose a harmless $T$ of size $b+1$; adding
the term $\ind\{T\subseteq A\}$ preserves all allowed responses but changes the
response at $T$. Removing subsumed terms leaves an antichain of rank at most $r$.
These alternatives prove necessity.
Conversely, under~\eqref{eq:radius}, all true edges and all negative candidates in
$\VV_r$ are queryable. Their labels determine all positive and negative members of
$\VV_r$. Every positive mask of a rank-at-most-$r$ monotone function contains a
positive candidate in $\VV_r$, so the entire function is determined.
\end{proof}
The usual condition $b\ge r$ is sufficient uniformly over the class, but is not
necessary for every instance. For six disjoint failure pairs, actual rank is two
and $\alpha=6$. Thus $b=6$ allows a rank-unrestricted certificate on twelve units;
$b<6$ does not. The largest observed failure order alone does not determine this
threshold.

\begin{theorem}[Exact identifying-certificate complexity]\label{thm:certificate}
When~\eqref{eq:radius} holds, the smallest truthful transcript uniquely identifying
$\HH$ has
\begin{equation}
 \Cert_{b,r}(\HH)=C_{b,r}(\HH)+s
 \label{eq:certificate}
\end{equation}
nonempty queries. Otherwise no allowed identifying transcript exists.
\end{theorem}
\begin{proof}
If some $T\in\II_r(\HH)$ is contained in no observed negative query, adding $T$ as a
positive term preserves the transcript while changing the function. Therefore the
negative queries form a cover and number at least $C_{b,r}(\HH)$. For each $e\in\HH$,
some positive query must contain $e$ and no other edge, or removing $e$ preserves
the transcript. A query can isolate at most one edge, so there are at least $s$
positive queries. This proves the lower bound.
For the upper bound, query all true minimal edges and an optimal negative cover.
An alternative consistent function must keep every true edge positive, and every
member of $\II_r(\HH)$ negative. Hence it agrees on $\VV_r$ and everywhere.
The necessity part of Theorem~\ref{thm:radius} rules out certificates otherwise.
\end{proof}
The negative-cover lower bound is classical; we state its exact certificate form
for the audit domain. An optimal certificate depends on the true target. Equation~
\eqref{eq:certificate} is not an algorithm that knows which optimal queries to ask.

\subsection{An independently checkable residual}
For discovered, verified edges $\KK$ and observed negative masks $\NN$, maintain
\begin{equation}
 \UU(\KK,\NN)=\{T\in\VV_r:
       \nexists e\in\KK\text{ with }e\subseteq T,\quad
       \nexists A\in\NN\text{ with }T\subseteq A\}.
 \label{eq:residual}
\end{equation}
An edge has a positive witness and negative witnesses for each single-element
deletion; a negative superset of a deletion is sufficient under monotonicity.
Also require every queried positive to contain a verified edge. Then $\UU=\varnothing$
is a complete certificate in the declared class. If $T\in\UU$, $B_{\KK}$ and
$B_{\KK}\vee\ind\{T\subseteq A\}$ both fit the transcript but disagree at $T$.
The alternative can be minimized without exceeding rank $r$.

Our checker validates query-domain membership, absence of positive-subset/negative-
superset contradictions, minimality witnesses, explained positives, the rank of the
reported edges, and an empty residual. It does not establish global monotonicity or
validate an unobserved rank bound. Detecting no contradiction is not a proof of
either assumption.

\section{Coverage-directed discovery with a global query bound}
For an allowed mask avoiding all found edges, define
\begin{equation}
 g(A;\UU)=|\{T\in\UU:T\subseteq A\}|.
\end{equation}
Our scheduler chooses a mask maximizing this score. A negative answer removes its
candidate subsets. A positive answer is shrunk to a minimal positive by restoring
one component at a time, retaining each restoration that leaves the mask positive.
All negative answers during shrinking also update the residual. The new edge removes
all of its candidate supersets. The process ends only when no residual candidates
remain, or reports a budget/domain/measurement limitation.

\begin{center}
\begin{minipage}{0.94\linewidth}
\hrule\vspace{5pt}
\textbf{Coverage-directed audit}\quad Inputs: $n,r,b$, oracle, query budget.\\
Initialize $\KK=\NN=\varnothing$, $\UU=\VV_r$.\\
While $\UU\ne\varnothing$:\\
\hspace*{1em}Choose $A\in\arg\max\{g(A;\UU):|A|\le b,\ \forall e\in\KK,\ e\nsubseteq A\}$.\\
\hspace*{1em}If the maximum is zero, report unresolved candidates and stop.\\
\hspace*{1em}Query $B(A)$, or abstain if its label cannot be resolved.\\
\hspace*{1em}If negative, add $A$ to $\NN$.\\
\hspace*{1em}If positive, deletion-shrink $A$ to $e$; add all shrink negatives to $\NN$; add $e$ to $\KK$.\\
\hspace*{1em}Recompute $\UU$ using~\eqref{eq:residual}.\\
Return $\KK$ with the labeled transcript and conditional certificate.\vspace{5pt}\hrule
\end{minipage}
\end{center}

\begin{theorem}[Global certificate-relative query bound]\label{thm:greedy}
Assume exact labels, exact score maximization, and~\eqref{eq:radius}. Let
\begin{equation}
 d=\min\left\{M,\sum_{j=1}^{\min(r,b)}\binom bj\right\},\qquad
 h_d=\sum_{j=1}^d\frac1j,\quad h_0=0.
\end{equation}
Coverage-directed auditing recovers $\HH$ and certifies completeness using at most
\begin{equation}
 Q\le h_d C_{b,r}(\HH)+(b+1)s.
 \label{eq:mainbound}
\end{equation}
In particular, $Q/\Cert_{b,r}(\HH)\le\max\{h_d,b+1\}$.
\end{theorem}
\begin{proof}
\emph{Correctness and progress.} A positive outer mask contains no previously found
edge. Deletion shrinking therefore finds a new true edge. If a deletion tested
negative before subsequent restorations, its final counterpart is a subset of that
negative mask and remains negative. Hence the returned edge is minimal.
If a residual candidate is positive, it contains an undiscovered, queryable edge;
querying that edge is feasible with positive score. If a residual candidate is
negative, some mask in an optimal negative cover contains it; this mask avoids all
found edges and also has positive score. Thus the algorithm progresses until the
residual is empty. Each original positive pool contains its extracted edge, so the
transcript certificate is valid.

\emph{One charge across all discovery phases.} Fix an optimal negative cover $\FF$
and assign each member of $\II_r(\HH)$ to one covering mask, yielding disjoint blocks
$P_F$ with $|P_F|\le d$. True-negative candidates disappear only through negative
queries, never by discovery of a true edge. Every $F\in\FF$ remains feasible at
all times. On an outer negative query, suppose the chosen mask covers $k$ residual
candidates. All $k$ are true negatives. Greedy optimality implies that $k$ is at
least the number $q$ of unresolved elements in any fixed block $P_F$.
Charge $1/k$ to each of the $k$ eliminated candidates. When $t$ elements of a block
with $q$ elements remaining are removed, its charge is at most
\begin{equation}
 \frac tq\le\sum_{j=q-t+1}^{q}\frac1j.
\end{equation}
Shrink-generated negative queries can delete additional block elements at zero
charge. Therefore, over the entire run, each block accrues at most $h_{|P_F|}\le h_d$.
The total number of outer negative queries is at most $C_{b,r}(\HH)h_d$.
There are at most $s$ outer positive queries, each followed by at most $b$ deletion
queries. Adding these costs proves~\eqref{eq:mainbound}. Caching only reduces calls.
\end{proof}
This uses the classical harmonic set-cover charge \citep{chvatal1979}, but the
cover's negative members are unknown to the algorithm and positive discoveries
change the feasible domain. Fixing the true negative cover throughout the proof
avoids paying its cost again after each discovery. The argument makes no overlap
or maximum-degree assumption.

\begin{corollary}[Approximate selection and alternative extraction]
If every chosen mask attains at least $1/a$ of the maximum score, the first term in
\eqref{eq:mainbound} becomes $a h_d C_{b,r}(\HH)$. If a correct minimal-edge extractor
uses at most $L$ additional queries per positive pool, the second term becomes
$(L+1)s$.
\end{corollary}
The first claim follows by replacing $k\ge q$ with $k\ge q/a$ in the same charge.
The second follows by replacing the extraction accounting. A practical heuristic
is not automatically assigned an approximation factor. Our implementation enumerates
$\sum_{j\le r}\binom nj$ candidates and $\sum_{j\le b}\binom nj$ masks; we claim a query
bound, not a polynomial-time method for unrestricted $n,r,b$. Existing efficient
hidden-hypergraph algorithms remain relevant alternatives \citep{angluin2006,abasi2014,chen2026}.

\section{The cost of excluding higher-order alternatives}
The declared order bound can matter far more than the actual order of the target.
Let $\HH_m$ consist of $m\ge2$ disjoint pairs on $n=2m$ components. Every component
participates and the true rank is two. Its maximal negative masks choose exactly
one component from each pair and have size $m$.

\begin{theorem}[Rank uncertainty and binary covering arrays]\label{thm:rankgap}
For $b\ge m$, $2\le r\le2m$, and $t=\min(r,m)$, $C_{b,r}(\HH_m)$ equals the smallest
number of rows in a binary strength-$t$ covering array with $m$ columns. In particular,
\begin{equation}
 C_{b,2m}(\HH_m)=2^m,\qquad
 C_{b,2}(\HH_m)\le2+2\lceil\log_2m\rceil,\qquad C_{b,r}(\HH_m)\ge2^t.
 \label{eq:rankgap}
\end{equation}
\end{theorem}
\begin{proof}
A harmless candidate chooses at most one vertex per pair. Extend each negative
query to a maximal one, and represent it by a binary row specifying which vertex
of each pair it chooses. Covering every candidate of at most $r$ vertices is exactly
the requirement that every assignment on any $t$ columns occurs in a row. Smaller
assignments extend to $t$ columns. This proves the equivalence.
For $t=m$, each binary string requires its own row. For $t=2$, label the columns
by distinct $\lceil\log_2m\rceil$-bit strings. Use the all-zero and all-one rows,
and each bit-position row with its complement. Constant rows cover the equal
assignments; distinct column labels differ at a bit, whose row and complement
cover both unequal assignments. Finally, any fixed $t$ columns require their
$2^t$ assignments in different rows.
\end{proof}
Covering arrays are classical \citep{colbourn2004}; the reduction makes the cost of
the audit's rank claim explicit. By Theorem~\ref{thm:certificate}, adding $m$ positive
pair witnesses gives a certificate of size at most $m+2+2\lceil\log_2m\rceil$ under a
trusted rank-two bound. Without that bound, the same target needs $2^m+m$ queries
in any exact identifying certificate. An externally known tight edge count would
exclude the extra-edge alternatives and invalidate this particular comparison.

For $b<m$ and $t=\min(r,m)\le b$, a negative mask covers at most $\binom bt$ of the
$2^t\binom mt$ oriented $t$-subsets. Thus
\begin{equation}
 C_{b,r}(\HH_m)\ge\left\lceil\frac{2^t\binom mt}{\binom bt}\right\rceil.
\end{equation}
For $b<t$, the negative cover is impossible. Positive edges must also be queryable.
This combines an intervention-size limit with an order-of-exclusion limit.

\section{Measured losses and abstention}\label{sec:noise}
A threshold decision from finitely many inputs can corrupt both elimination and
minimality witnesses. We use a deliberately conservative standard concentration
wrapper, rather than introduce a new statistical method. For distinct query $i$,
draw fresh independent losses in $[0,1]$. At sample count $u$, let
\begin{equation}
 \epsilon_{i,u}=\sqrt{\frac{\log(\pi^4 i^2u^2/(18\delta))}{2u}}.
\end{equation}
Return positive when $\widehat D-\epsilon_{i,u}>\tau$, negative when
$\widehat D+\epsilon_{i,u}\le\tau$, and otherwise gather more samples or abstain.
Hoeffding's inequality gives conditional failure probability at most
$36\delta/(\pi^4i^2u^2)$. Summing over $i,u\ge1$ bounds the probability of \emph{any}
incorrect accepted label by $\delta$, including adaptive mask choices. On the
complement event the deterministic results apply unchanged. This justification
requires fresh samples for each query, not unrestricted reuse of a fixed prompt set.
A margin $|D(A)-\tau|\ge\gamma$ permits termination once $2\epsilon_{i,u}<\gamma$.
A sample cap can instead yield a legitimate unresolved measurement. We separately
report mask requests and total input samples.

\section{Experiments}
\subsection{Protocols and comparison objectives}
All runs use CPU computation. Synthetic oracles are exact and neural intervention
tables are exhaustively enumerated. Truth is used only for evaluation and optimal-
certificate calculations; discovery and stopping receive only queried labels.
The release contains seeds, raw rows, an example transcript, model weights for seed
zero, exact optimizer outputs, tests, and scripts regenerating the tables and figures.
The 2,880-audit synthetic sweep plus neural and noise suites took 251.46 seconds in
our recorded container run. The additional rank sweep and independent proof checks
are separate runs; these are not laptop timing guarantees.

\textbf{Baselines.} Random~+~deletion shares the residual ledger and minimal-set
extractor with our method, but samples uniformly among largest-size informative
masks avoiding known edges. This isolates the outer scheduling decision. SIGHT-style
uses binary-prefix extraction followed by deletion verification. RC-style uses
random shrinking (twenty trials per reduction), then increasing-order subset tests;
a repeatedly unproductive pool falls back to deletion after three failures. Both
use the shared full-recovery scheduler, rather than the original one-edge sampling
protocols of \citet{clarfeld2019}. Increasing-order enumeration tests candidate
orders from smallest upward and skips already explained or excluded candidates.
The pairs-only control refuses to claim completeness while higher-order candidates
remain. We do not present these adaptations as tuned reproductions of all published
implementations, nor compare their query counts to transformer-specific CoAx scores.

\subsection{Independent checks of the mathematical statements}
We enumerate all 167 monotone functions on four components with harmless empty mask.
Across all declared ranks and intervention ceilings from one to four, 1,848 cases
satisfy the predicted identifiability characterization exactly. For each of the
1,254 identifiable cases we solve a separate integer program selecting a minimum
transcript that distinguishes the target from \emph{every} alternative hypothesis.
Its optimum equals $C+s$, computed by an independent negative-cover program.
Coverage-directed auditing recovers each target, passes the transcript checker, and
satisfies~\eqref{eq:mainbound}. An additional 332 rank-three cases check the same
bounds, and 42 unit tests cover invariants, budgets, noise abstention, and deliberate
assumption failures. Finite verification supplements, and does not replace, the proofs.

\subsection{Synthetic query efficiency}
We use twelve units, declared rank three, ceilings $b\in\{3,4,6,8\}$, and twenty seeds.
The six fixed families are empty; four disjoint triples; eight triples with a shared
pair; six randomly selected triples; a mixed-order antichain; and all twenty triples
within a six-unit module. Seeds permute component labels and algorithm tie-breaking;
the random-triple family also changes structure. Each method is evaluated on the
same oracle for a given family, ceiling, and seed.

\begin{table}[t]
\centering\small
\begin{tabular}{lrrrr}
\toprule
Method & $b=3$ & $b=4$ & $b=6$ & $b=8$ \\
\midrule
Coverage-directed & $221.2 \pm 0.0$ & $92.9 \pm 0.2$ & $64.6 \pm 0.1$ & $69.9 \pm 0.2$ \\
Random + deletion & $221.2 \pm 0.0$ & $130.3 \pm 0.3$ & $80.0 \pm 0.3$ & $76.0 \pm 0.2$ \\
SIGHT-style & $224.9 \pm 0.1$ & $141.0 \pm 0.5$ & $92.3 \pm 0.5$ & $85.0 \pm 0.5$ \\
RC-style & $228.7 \pm 0.1$ & $142.0 \pm 0.3$ & $116.6 \pm 1.1$ & $112.6 \pm 1.5$ \\
Increasing order & $284.2 \pm 0.0$ & $284.2 \pm 0.0$ & $284.2 \pm 0.0$ & $284.2 \pm 0.0$ \\
\bottomrule
\end{tabular}

\caption{Queries to a checked completeness certificate. Entries are means $\pm$ standard
errors over twenty seed-wise averages of the six fixed families, not population
uncertainty over arbitrary circuits. All five methods recover and certify all 480
instances. The pairs-only control never certifies the rank-three class.}
\label{tab:queries}
\end{table}

Table~\ref{tab:queries} shows the main query comparison. At $b=6$, coverage-directed
selection uses 64.6 queries on average, compared with 80.0 for matched random
selection. Across all 480 paired instances, it uses fewer queries on 362, the same
number on 108, and more on 10. The reduction in the overall mean is 11.6\%.
These are descriptive results on the fixed benchmark, not a universal ordering of
algorithms. Increasing $b$ does not monotonically improve the greedy scheduler:
the dense-module family costs 137.2 queries at $b=4$ and 181.0 at $b=8$, where positive
pool extraction becomes expensive. A larger feasible domain does not force an
algorithm to choose cost-effective pools.

\begin{figure}[t]
\centering\includegraphics[width=.78\linewidth]{../figures/query_cost.pdf}
\caption{Synthetic certificate cost versus intervention ceiling. Error bars follow
Table~\ref{tab:queries}. Query cost, not optimizer runtime, is the evaluation target.}
\label{fig:queries}
\end{figure}

\subsection{Holding the network fixed while changing the claim}
For six disjoint pairs on twelve units, fix $b=6$ and vary only the declared rank
$r$. We run four auditors for twenty seeds at each of six ranks (480 runs).
Integer optimization proves the negative-cover values in Table~\ref{tab:rank}.
The true mechanism does not change. Coverage-directed discovery finds all six pairs
at approximately 44 queries for every rank, but certification grows from 46.25 to
106 queries as higher-order alternatives are admitted.

\begin{table}[t]
\centering\small\begin{tabular}{rrrrr}
\toprule
Declared $r$ & Min. negative cover & Discovery & Certification & Min. certificate \\
\midrule
2 & 6 & 43.55 & $46.2 \pm 0.2$ & 12 \\
3 & 12 & 43.60 & $54.2 \pm 0.2$ & 18 \\
4 & 21 & 43.75 & $68.3 \pm 0.3$ & 27 \\
5 & 32 & 43.75 & $86.2 \pm 0.9$ & 38 \\
6 & 64 & 43.75 & $106.0 \pm 0.0$ & 70 \\
12 & 64 & 43.75 & $106.0 \pm 0.0$ & 70 \\
\bottomrule
\end{tabular}

\caption{The same six failure pairs under different rank assumptions. Discovery is
an evaluator-only stopping time; certification is the actual audit stopping time.
Minimum certificate is an oracle-dependent lower benchmark, not an implementable
policy with advance knowledge of the target. Query uncertainty is over twenty seeds.}
\label{tab:rank}
\end{table}

\begin{figure}[t]
\centering\includegraphics[width=.78\linewidth]{../figures/rank_discovery_gap.pdf}
\caption{Discovering all true pairs and excluding unobserved higher-order failures
have different costs. The curves use the same frozen oracle at every rank.}
\label{fig:rank}
\end{figure}

A further 100 runs hold $r=12$ and vary $b$ from two to six. At $b=2,3,4,5$, the
residual contains respectively 656, 496, 256, and 64 candidate sets after all allowed
inferences have been exhausted. At $b=6$ it is empty. These counts equal
$\sum_{k=b+1}^{6}2^k\binom6k$ and verify the sharp radius $b_*=6$.
Separate random-pool experiments match the elementary single-hidden-triple detection
formula given in Appendix~\ref{app:locality}; they are sanity checks, not evidence
that coverage-directed scheduling is optimal.

\subsection{Small trained networks}
We train twenty ReLU classifiers: ten seeds each with nonnegative and signed final
readouts. Inputs are all 256 vectors in $\{-1,1\}^8$; the target is the nonnegative
sum of the first four bits. Hidden width is 32, grouped into eight fixed intervention
units. Adam uses learning rate 0.02 for 1,200 full-domain steps with independent
20\% group dropout. We freeze each model and enumerate all 256 masks under zero and
mean replacement, giving forty model/policy conditions and 2,621,440 mask--input
pairs. The discrepancy is prediction disagreement and $\tau=0.10$.
All models attain 100\% accuracy on this finite training/evaluation universe; we make
no held-out distributional-generalization claim.

The positive-readout, zero-replacement setting has a structural monotonicity proof
(Appendix~\ref{app:neural}). Exhaustive adjacent-mask checks also happen to find no
violations in the other sampled settings; this is an empirical property of these
models, not a guarantee for signed readouts or mean replacement. Actual minimal
failure sets are obtained from each frozen model's response table, never assumed
from the task or architecture. All auditors use $r=b=n=8$, without access to the true
failure order. The 160 audits all recover their exact tables and pass the certificate
checker. Precomputing the full tables is a validation cost, not counted as queries
available to discovery.

\begin{table}[t]
\centering\small\setlength{\tabcolsep}{4pt}\begin{tabular}{lrrrrrr}
\toprule
Readout / patch & Edges & Max. rank & Coverage & Random & SIGHT & Increasing \\
\midrule
Positive / zero & 7.8 & 7--7 & 38.5 & 38.4 & 99.4 & 253.3 \\
Positive / mean & 25.0 & 2--2 & 71.4 & 71.3 & 54.7 & 32.5 \\
Signed / zero & 1.2 & 7--8 & 11.5 & 11.1 & 25.2 & 254.7 \\
Signed / mean & 42.3 & 5--5 & 156.3 & 154.2 & 203.1 & 182.3 \\
\bottomrule
\end{tabular}

\caption{Means across ten seeds per model/policy condition. ``Max. rank'' is the
range of each model's largest minimal failure-set size. Last four columns count
queries to a certificate. All methods are exact on these forty conditions.}
\label{tab:neural}
\end{table}

The results delimit the scheduling advantage. On positive-readout zero-ablation,
coverage-directed and random selection use 38.5 and 38.4 queries, respectively;
increasing-order enumeration uses 253.3. Under positive-readout mean replacement,
increasing order instead uses 32.5 queries, versus 71.4 for coverage-directed
selection. Thus the neural experiments validate conditional correctness and the
importance of intervention semantics, not uniform superiority of our scheduler.

\subsection{Noise and explicit assumption failures}
For an eight-unit two-triple oracle, set each query's Bernoulli loss mean to
$0.5+\gamma$ if harmful and $0.5-\gamma$ otherwise, with $\gamma\in\{0.1,0.2\}$.
At each margin we compare twenty confidence-wrapped runs ($\delta=0.05$) with twenty
runs that threshold a fixed batch of 32 samples. The confidence wrapper resolves
labels at geometrically increasing sample counts, capped at 262,144 per query.

\begin{table}[t]
\centering\small\setlength{\tabcolsep}{4pt}\begin{tabular}{rlrrrrr}
\toprule
Margin & Labels & Masks & Input samples & Exact & Checked cert. & False cert. \\
\midrule
0.1 & Confidence & 26.6 & 45,030 & 20/20 & 20/20 & 0/20 \\
0.1 & Fixed 32 & 33.7 & 1,078 & 0/20 & 4/20 & 4/20 \\
0.2 & Confidence & 26.6 & 10,138 & 20/20 & 20/20 & 0/20 \\
0.2 & Fixed 32 & 26.3 & 842 & 16/20 & 19/20 & 3/20 \\
\bottomrule
\end{tabular}

\caption{Noisy auditing. ``Checked cert.'' passes the independent transcript checker;
``False cert.'' passes that checker but identifies the wrong true oracle. The checked
false count excludes contradictions already caught by the checker. Mask and sample
columns are means; success columns count twenty runs.}
\label{tab:noise}
\end{table}

All forty confidence-wrapped runs are exact with no erroneous accepted labels, at
substantially greater sample cost. The fixed-batch baseline produces seven false
\emph{checked} certificates across forty runs; an additional five incorrect stopping
claims are caught by the checker. Forty successful confidence runs do not themselves
establish the statistical guarantee; its justification is the concentration argument.

We also construct assumption failures from a frozen trained network. Add two constant
hidden units with equal and opposite final contributions $+a,-a$, where
$a>\max_x|f_{\rm logit}(x)|$. Clean logits are unchanged. Considering these two units
as the intervention universe yields $B(\{1\})=B(\{2\})=1$ but $B(\{1,2\})=0$.
The monotone audit can query the pair, see a negative response, and issue a
\emph{conditional} empty-hypergraph certificate that is wrong for this non-monotone
network. The transcript has no contradiction. Likewise, declaring rank two and
allowing only pairs on a trained higher-order model yields an unsupported empty
explanation. Allowing a larger query exposes a rank violation and triggers an
explicit warning. These are constructive controls, not claims that such patterns
occur naturally at a measured frequency.

\section{Related work and limitations}
\textbf{Prior mathematics.} Hidden-hypergraph and monotone-DNF learning already
provide identification algorithms, lower bounds, and independent-cover constructions
\citep{angluin2006,abasi2014,chen2026}. Our negative certificate and covering-array
objects are specializations of classical constructions, and harmonic charging is
standard \citep{chvatal1979,colbourn2004}. The proposed technical emphasis is the
single certificate-relative charge across adaptive discoveries under bounded
interventions, together with a rank-sensitive interpretation of completeness.
We do not claim improved unrestricted worst-case hypergraph-learning complexity.

\textbf{Interpretability and group testing.} Formal mechanistic interpretability
uses verification to provide circuit guarantees with different minimality and
robustness objectives \citep{hadad2026}. CoAx studies conditional backup recovery
and transformer self-repair \citep{gong2026}; its seeded component-level objective
differs from recovery of an entire higher-order failure hypergraph. SIGHT and Random
Chemistry motivate our group-extraction baselines \citep{clarfeld2019}. Sparse
M\"obius methods recover polynomial interactions under corresponding sparsity
assumptions \citep{kang2024}; sparsity of minimal positive terms is not the same as
sparsity of a Boolean polynomial after inclusion--exclusion. We have not benchmarked
these external implementations on their native tasks.

\textbf{What is not certified.} The model class may be false: signed effects,
restoration under additional ablations, context-specific thresholds, or an incorrect
rank bound can invalidate inference. A certificate concerns the fixed $P$, policy,
components, and threshold in~\eqref{eq:behavior}; changing them changes the question.
It does not recover a unique causal graph, assign semantic equivalence to redundant
units, or certify behavior on unobserved distributions. The noise guarantee requires
fresh sampling and makes no claim of sample-optimality. Finally, exact mask selection
and residual enumeration can be exponential; scaling computation without losing the
query guarantee is an unresolved engineering and theoretical extension of this work.

\section{Conclusion}
Finding dependencies and certifying their completeness require different evidence.
We characterize the intervention radius and certificate size, prove a
certificate-relative adaptive bound, and isolate the exponential cost of excluding
higher-order alternatives. The experiments show both gains and failures of greedy
scheduling. A defensible explanation must retain its intervention scope,
structural assumptions, labeled evidence, and unresolved alternatives.

\begin{thebibliography}{9}
\bibitem[Abasi et~al.(2014)Abasi, Bshouty, and Mazzawi]{abasi2014}
Hasan Abasi, Nader H. Bshouty, and Hanna Mazzawi.
\newblock On Exact Learning Monotone DNF from Membership Queries.
\newblock arXiv:1405.0792, 2014. \url{https://arxiv.org/abs/1405.0792}.

\bibitem[Angluin and Chen(2006)]{angluin2006}
Dana Angluin and Jiang Chen.
\newblock Learning a Hidden Hypergraph.
\newblock \emph{Journal of Machine Learning Research}, 7:2215--2236, 2006.
\newblock \url{https://www.jmlr.org/papers/v7/angluin06a.html}.

\bibitem[Chen(2026)]{chen2026}
Ting Chen.
\newblock More Efficient Algorithms for Searching for Several Edges in a Hypergraph.
\newblock \emph{Indian Journal of Pure and Applied Mathematics}, 57:425--433, 2026.
\newblock First published online in 2024. doi:10.1007/s13226-024-00561-z.

\bibitem[Chv\'atal(1979)]{chvatal1979}
V\'aclav Chv\'atal.
\newblock A Greedy Heuristic for the Set-Covering Problem.
\newblock \emph{Mathematics of Operations Research}, 4(3):233--235, 1979.
\newblock doi:10.1287/moor.4.3.233.

\bibitem[Clarfeld and Eppstein(2019)]{clarfeld2019}
Laurence A. Clarfeld and Margaret J. Eppstein.
\newblock Group-Testing on Hypergraphs with Variable-Cost Tests: A Power Systems Case Study.
\newblock arXiv:1909.04513, 2019. \url{https://arxiv.org/abs/1909.04513}.

\bibitem[Colbourn(2004)]{colbourn2004}
Charles J. Colbourn.
\newblock Combinatorial Aspects of Covering Arrays.
\newblock \emph{Le Matematiche}, 59(1--2):125--172, 2004.
\newblock \url{https://lematematiche.dmi.unict.it/index.php/lematematiche/article/view/166}.

\bibitem[Gong et~al.(2026)Gong, Zeng, Yuen, and Lim]{gong2026}
Zhiren Gong, Zihao Zeng, Chau Yuen, and Wei Yang Bryan Lim.
\newblock Conditional Co-Ablation: Recovering Self-Repair Backups in Transformer Circuits.
\newblock arXiv:2607.01940, 2026. \url{https://arxiv.org/abs/2607.01940}.

\bibitem[Hadad et~al.(2026)Hadad, Katz, and Bassan]{hadad2026}
Itamar Hadad, Guy Katz, and Shahaf Bassan.
\newblock Formal Mechanistic Interpretability: Automated Circuit Discovery with Provable Guarantees.
\newblock arXiv:2602.16823, 2026. \url{https://arxiv.org/abs/2602.16823}.

\bibitem[Kang et~al.(2024)Kang, Erginbas, Butler, Pedarsani, and Ramchandran]{kang2024}
Justin S. Kang, Yigit E. Erginbas, Landon Butler, Ramtin Pedarsani, and Kannan Ramchandran.
\newblock Learning to Understand: Identifying Interactions via the M\"obius Transform.
\newblock arXiv:2402.02631, 2024. \url{https://arxiv.org/abs/2402.02631}.
\end{thebibliography}


\appendix
\section{Elementary locality calculation}\label{app:locality}
Let there be either no failure edge or one hidden $k$-set $T$, with
$B_T(A)=\ind\{T\subseteq A\}$. If $b<k$, every allowed query returns zero in both
cases. If $b\ge k$, each allowed mask covers at most $\binom bk$ of the $\binom nk$
possible targets. Thus the all-negative branch of any exact audit certifying absence
requires at least $\lceil\binom nk/\binom bk\rceil$ queries. Adaptivity does not
change this covering argument along that branch. This is only a counting lower
bound; an optimal covering design may require more masks.

An independently sampled uniform size-$b$ mask contains a fixed target with
probability $p=\binom bk/\binom nk$, so the probability of discovery after $Q$ masks
is $1-(1-p)^Q$. We simulate actual hypergeometric overlaps for $n=32,k=3$, seven
ceilings, five budgets, and 2,000 trials per cell. The exact limits at $b=4$ and $b=8$
are 1,240 and 89 queries, respectively, for the counting lower bound.

\begin{figure}[h]
\centering\includegraphics[width=.75\linewidth]{../figures/locality.pdf}
\caption{Single-hidden-triple sanity check. Lines give the analytic random-pool
detection probability; markers are empirical frequencies. These curves describe
random sampling, not a zero-error completeness certificate.}
\end{figure}

\section{Why the constrained neural model is monotone}\label{app:neural}
Write the unablated logit as $z(x)=c+\sum_{j=1}^{n}v_j(x)$, where each group
contribution $v_j(x)$ is nonnegative. Zero-ablation gives
$z_A(x)=c+\sum_{j\notin A}v_j(x)$. Consequently $A\subseteq A'$ implies
$z_{A'}(x)\le z_A(x)$. If the unablated prediction is negative, it stays negative;
if positive, nested ablations can change it only once, from positive to negative.
Thus $\ind\{\widehat y_A(x)\ne\widehat y_\varnothing(x)\}$ is monotone for every
input. Expectation over inputs and thresholding preserve monotonicity.

Mean replacement instead removes $v_j(x)-\E[v_j(x)]$, which need not be nonnegative.
Signed readouts also lose the sign argument. Our finite-table checks for these
settings are additional measurements, not consequences of the architecture.
The cancellation control is representable by two constant ReLU units with final
weights $+a$ and $-a$, so it is a neural counterexample with unchanged clean logits,
not merely an arbitrary Boolean oracle.

\section{Reproducibility, baseline details, and artifact scope}
The experiment package uses Python, NumPy, SciPy, pandas, matplotlib, and CPU PyTorch.
The recorded environment uses Python 3.13.5, NumPy 2.3.5, SciPy 1.17.0, pandas 2.2.3,
matplotlib 3.10.8, and PyTorch 2.10.0+cpu. Training uses one PyTorch thread and
deterministic algorithms. Algorithms operate on integer bitmasks; bitset candidate
coverage makes exhaustive twelve-component experiments inexpensive. Model weights
and discrepancy tables for seed zero are released as JSON; all twenty models are
reproducible from their training seeds, but checkpoints for the other seeds are not
included in this release.

The increasing-order baseline is not forced to evaluate every possible mask: it
shares negative-subset and positive-superset inference with the other methods. Random,
SIGHT-style, and RC-style schedulers also avoid masks containing known edges and do
not waste a test on a mask with zero residual coverage. These are stronger common-
ledger controls than unconditioned random masks. SIGHT-style uses binary-prefix
extraction and an extra deletion verification; this may cost more than an optimized
extractor when the minimal edge is large. RC-style may fail to shrink a positive
pool within its twenty-trial limit; it retries rather than silently discarding that
pool, and eventually uses deletion. The main theorem concerns coverage-directed
selection with deletion, not the query complexity of these adaptations.

Negative-cover integer programs are restricted to maximal allowed harmless masks,
which dominate their harmless subsets for uniform query costs. Solver outputs are
called optima only when the mixed-integer solver reports optimality. The independent
minimum-transcript program instead has one constraint per alternative hypothesis:
at least one selected query must distinguish that alternative from the target.
The agreement of these independently formulated optimizations is reported only for
the small cases that were actually solved.

\begin{table}[h]
\centering\small\begin{tabular}{rrrl}
\toprule
Ceiling $b$ & Queries & Residual candidates & Output \\
\midrule
2 & 78.0 & 656 & Locality limited \\
3 & 183.9 & 496 & Locality limited \\
4 & 269.9 & 256 & Locality limited \\
5 & 228.0 & 64 & Locality limited \\
6 & 106.0 & 0 & Complete \\
\bottomrule
\end{tabular}

\caption{Intervention-radius sweep with six disjoint pairs, declared rank twelve,
and twenty tie-breaking seeds per ceiling. All six true pairs are eventually found
at every ceiling, but only $b=6$ permits full identification.}
\end{table}

\textbf{Reproduction commands.} From the repository root, install the package with
its experiment extras, then run \texttt{make test} and \texttt{make experiments}.
\texttt{make paper} regenerates figures and tables from the released raw records
and compiles this manuscript. The packaged raw-record archive can be expanded with
\texttt{python experiments/unpack\_results.py}. No LLM API, downloaded model weights,
or GPU is used. The full ground-truth enumeration belongs to validation; its cost
must not be confused with the oracle queries revealed to any auditing algorithm.

\textbf{Release status.} This is a reproducible research manuscript, not a claim of
acceptance or independent verification. Historical priority of the proposed
certificate-relative adaptive bound still requires a dedicated external comparison
with the exact-learning literature. Further practical evaluation should use an
external circuit benchmark and an optimizer that does not enumerate the whole
intervention space. The present conclusions are restricted to the proofs and
finite experiments stated here.
\end{document}
